\chapter{背景知识}

1. 什么是 Dify？

**Dify** 是一个**开源的大语言模型（LLM）应用开发平台**，旨在帮助开发者快速构建和部署基于 AI 的应用。它提供了一个直观的界面，集成了 AI 工作流、RAG（检索增强生成）、智能体（Agent）功能和模型管理等工具。简单来说，Dify 就像一个“搭积木”的平台，用户无需深入编码就能创建 AI 应用，例如智能客服或知识库问答系统。

 2. 什么是 Ollama？
 
**Ollama** 是一个**轻量级的本地推理框架**，用于在本地机器上部署和运行大语言模型（LLM）。它支持多种模型（如 LLaMA、Mistral、DeepSeek 等），用户只需通过简单命令即可下载和运行这些模型。Ollama 的优势在于本地化部署，不依赖云服务，保证数据隐私，同时提供 API 接口供其他应用调用。

 3. 什么是 Embedding？
 
**Embedding（嵌入）** 是一种**将文本、图像等数据转化为向量表示的技术**。在 AI 领域，Embedding 模型将复杂的自然语言或数据映射成数字向量，这些向量保留了语义信息。例如，“猫”和“狗”的向量会比较接近，而“猫”和“汽车”的向量距离较远。Embedding 在知识检索、语义搜索和问答系统中起着关键作用。

4. 什么是 LightRAG？

**LightRAG** 是一个**轻量级的 RAG（检索增强生成）系统**，结合了知识图谱和向量检索技术，以提高数据检索的效率和准确性。它通过将文本数据转化为知识图谱，并利用向量检索快速定位相关信息，为大语言模型提供更精准的上下文。LightRAG 注重轻量化和高效性。

5. 什么是 GraphRAG？

**GraphRAG** 是一种**基于图的检索增强生成方法**，利用知识图谱增强大语言模型的检索能力。它通过构建文本数据的知识图谱，并利用图中的社区结构和层次关系，生成更全面和准确的回答。GraphRAG 特别适合处理复杂查询和需要深入理解数据集上下文的场景。

 6. 什么是 LazyGraphRAG？
 
**LazyGraphRAG** 是一个**类似于 GraphRAG 的系统**，但在构建和查询知识图谱时采用了更高效的策略，以减少计算资源需求。它可能通过延迟加载或按需计算图谱的某些部分来优化性能，从而在保持准确性的同时提升效率。

7. 什么是 MinerU？

**MinerU** 是一个**知识图谱构建工具**，能够从文本数据中自动提取实体和关系，并构建知识图谱。它通常与 RAG 系统结合使用，提供结构化的知识表示，从而提高检索和生成的准确性。

- 8. 什么是 DeepSeek？

**DeepSeek** 是一个**开源的大语言模型**，在 RAG 系统中常用作生成模型。它以高性能和低成本著称，特别在处理大规模文档和复杂查询时表现出色。DeepSeek 可通过 Ollama 在本地运行，提供隐私保护和低延迟的 AI 服务。

 它们之间的联系
这些工具和技术在 AI 和大语言模型生态系统中相互关联，协同工作以构建高效、私密且功能强大的应用：

- **Dify** 作为一个开发平台，可以集成 **Ollama** 提供的本地模型服务。用户可在 Dify 中配置 Ollama 的 API 接口，使用本地运行的模型（如 **DeepSeek**），从而构建隐私性强的 AI 应用。
- **Embedding** 是 RAG 系统的核心技术，用于将文本数据转化为向量表示，支持 **LightRAG**、**GraphRAG** 和 **LazyGraphRAG** 的语义搜索和检索。Ollama 可提供本地 Embedding 模型（如 nomic-embed-text）。
- **LightRAG**、**GraphRAG** 和 **LazyGraphRAG** 是 RAG 系统的变体，目标是通过检索外部知识增强模型生成能力。**LightRAG** 强调高效性，**GraphRAG** 利用知识图谱提升准确性和深度，**LazyGraphRAG** 则优化资源消耗。
- **MinerU** 与 **GraphRAG** 和 **LazyGraphRAG** 结合，从文本中提取实体和关系，构建知识图谱，为检索提供结构化支持。
- **DeepSeek** 作为生成模型，可通过 Ollama 在本地运行，并集成到 Dify 应用中，结合 RAG 系统提供准确回答。

总结
- **Dify**：应用开发平台。
- **Ollama**：本地模型运行框架。
- **Embedding**：文本向量化技术。
- **LightRAG**：轻量级 RAG 系统。
- **GraphRAG**：基于图的 RAG 系统。
- **LazyGraphRAG**：高效优化的 GraphRAG。
- **MinerU**：知识图谱构建工具。
- **DeepSeek**：高性能生成模型。

这些工具通过协同工作，支持从知识提取到检索和生成的全流程，适用于本地化、私密性强的 AI 应用场景。



\section{1Panel}
\url{https://1panel.cn/}
1Panel 提供了一个直观的 Web 界面，帮助用户轻松管理 Linux 服务器中的网站、文件、容器、数据库以及大型语言模型（LLMs）。

安装见官网。

需要注意的是：面板地址在安装设置好了，就不好修改。\\
可以用命令查阅：\verb|sudo 1pctl user-info|




\section{Docker}
Docker 是一个开源平台，用于自动化部署、运行和管理应用程序的容器化技术。它允许开发者将应用及其依赖打包到一个轻量、可移植的容器中，确保在不同环境（如开发、测试、生产）下运行一致。容器比虚拟机更高效，因为它们共享主机操作系统的内核，启动快、占用资源少。

主要特点：
\begin{itemize}
\item 容器化：隔离应用运行环境。
\item 可移植性：一次构建，任意环境运行。
\item 生态系统：支持 Docker Hub 共享镜像，简化部署。
\item 编排：常与 Kubernetes 等工具结合，管理大规模容器。
\end{itemize}
简单说，Docker 就像一个“标准化集装箱”，让应用打包和运行更简单高效

安装：\verb|sudo apt-get install docker-compose|

注意：更新Docker源，不然可能拉取失败。


\section{MongoDB}
MongoDB 是一个开源的 NoSQL 数据库，采用文档导向（document-oriented）存储方式。它以 JSON 格式的 BSON（Binary JSON）文档存储数据，灵活支持非结构化或半结构化数据，适合快速开发和处理大规模数据。

主要特点：
\begin{itemize}
\item 文档模型：数据以键值对形式存储，像 JSON 对象， schema 灵活。
\item 高性能：支持索引、内存映射，查询速度快。
\item 可扩展：通过分片和复制集实现水平扩展和高可用。
\item 易用：无需预定义表结构，适合动态变化的数据。
\end{itemize}
简单说，MongoDB 是一个高效、灵活的数据库，特别适合需要快速迭代或处理复杂数据的现代应用，比如 Web 应用、实时分析等。



\chapter{内网穿透、反向代理、异地组网}
如果学校网络不允许直接端口映射，但你有其他公网服务器，可以通过反向代理将请求转发到内网。

步骤：

在公网服务器上安装 Nginx 或 Apache。
配置反向代理，将公网服务器的某个域名或路径（如 example.com）转发到内网服务器。
使用内网穿透工具（如 Frp 或 SSH 隧道）将内网服务器连接到公网服务器。
优点：

灵活，可结合域名提供专业化访问。
适合有公网服务器的情况。
缺点：

需要额外的公网服务器，配置稍复杂。


\section{TLS}
TLS（传输层安全协议）是一种加密协议，用于在两个通信应用之间提供隐私和数据完整性。它主要用于保护Web通信（如HTTPS），确保客户端与服务器之间的数据传输是私密且经过身份验证的。TLS可以保证数据在传输过程中不被窃听或篡改。



\section{Nginx}
NGINX（发音为 "engine-X"）是一款开源的、高性能的HTTP和反向代理服务器，同时也提供了诸如IMAP/POP3代理服务器等功能。它最初由Igor Sysoev为俄罗斯大型网站如Rambler开发，并于2004年首次公开发布。NGINX以其高效率、稳定性和低资源消耗而著称。

作为一个HTTP服务器，NGINX可以用来直接服务静态内容（如HTML页面、样式表、脚本文件等），并且可以通过FastCGI、SCGI处理程序或者uWSGI协议来服务动态网页内容。此外，NGINX还广泛用于作为反向代理服务器，用于将客户端请求转发给后端服务器（例如运行着应用服务器、数据库服务器等），并将响应返回给客户端，这样可以提高安全性、性能和扩展性。

除了基本的服务器功能外，NGINX还支持负载均衡、HTTP缓存、SSL/TLS终止等功能，使其成为构建高效Web基础设施的关键组件。NGINX Plus是其商业版，提供额外的功能和官方支持。


NGINX是一款高性能的HTTP和反向代理服务器，同时支持IMAP/POP3代理功能。它因其高效率、稳定性和低资源消耗而广受欢迎。NGINX不仅能有效地服务静态内容，还能作为反向代理来处理动态内容请求、执行负载均衡、管理TLS终止等任务。

三者的关系
反向代理与NGINX：NGINX本身可以作为一个反向代理服务器使用。这意味着它可以接收来自互联网的客户端请求，然后根据配置将这些请求转发给一组后端服务器（例如应用服务器、数据库服务器等），并将响应返回给客户端。这样做的好处之一是可以增强安全性并优化性能。

TLS与NGINX：NGINX能够处理TLS加密和解密过程。这意味着它可以作为TLS终止点，接收来自客户端的加密请求，解密这些请求以便进一步处理或转发到适当的后端服务器，并对返回给客户端的响应进行加密。这不仅减轻了后端服务器的负担，还简化了证书管理。

反向代理与TLS：当结合使用时，反向代理（如NGINX）能够为后端服务器提供额外的安全层。通过在反向代理上终止TLS连接，所有进入的流量都会首先被加密和验证，从而保护内部网络不受外部威胁的影响。
综上所述，通过将NGINX用作反向代理并处理TLS加密，组织可以构建一个既安全又高效的Web服务架构。这种方式不仅可以提升应用的安全性，还可以通过负载均衡和缓存等技术显著改善性能。



\section{vnt}
\url{https://github.com/vnt-dev/vnt}



\section{frp}
\url{https://github.com/fatedier/frp/blob/dev/README_zh.md}

frp 是一个专注于内网穿透的高性能的反向代理应用，支持 TCP、UDP、HTTP、HTTPS 等多种协议，且支持 P2P 通信。可以将内网服务以安全、便捷的方式通过具有公网 IP 节点的中转暴露到公网。








\chapter{Overleaf}
Refs：
\begin{itemize}
\item 这里教程的主要参考：\url{https://www.cnblogs.com/Undefined443/p/18155463}

\item 官方教程：\url{https://github.com/overleaf/toolkit/blob/master/doc/quick-start-guide.md}

\item 官方说明文档：\url{https://github.com/overleaf/toolkit/tree/master/doc}

\item \url{https://sspai.com/post/83982}

\item \url{https://www.tnnidm.com/build-and-use-overleaf-server/index.html}
\end{itemize}

2017 年，Overleaf 和 ShareLaTeX 宣布合并，合并后的平台使用 Overleaf 作为名称。原先的 ShareLaTeX 代码库也改名为 Overleaf。不过目前官方提供的 Docker 映像仍叫 ShareLaTeX。因此下面有时会将 Overleaf 映像称为 ShareLaTeX 映像。

可以使用 overleaf-toolkit 来安装 overleaf （简单）或者使用 overleaf 的 docker compose 来安装 overleaf （复杂）

这里使用overleaf-toolkit。
注意：本教程只适用于 macOS / Linux 操作系统。如果需要在 Windows 上部署 Overleaf，请先安装 WSL，之后在 WSL 中部署 Overleaf。







\subsection{本地部署步骤}
0、安装 Docker。由于 toolkit 使用 docker-compose 作为容器编排器，没有 docker-compose 的同学还需先安装一下。\\
安装：\verb|sudo apt-get install docker-compose|

1、克隆 Overleaf Toolkit 仓库：\\
\verb|git clone https://github.com/overleaf/toolkit.git overleaf-toolkit|

2、进入 overleaf-toolkit 目录并生成配置文件：
\verb|cd overleaf-toolkit|\\
\verb|bin/init|

3、下载 Overleaf 映像并创建容器：
\verb|bin/up|

注意：
\begin{itemize}
\item Overleaf 脚本中使用的命令行工具均为 GNU 版本，而 macOS 原装命令行工具为 BSD 版本。因此 Mac 用户在启动时会报错。请使用 Homebrew 安装相应工具的 GNU 版本（如 coreutils），然后再尝试安装。

\item 要先改配置文件，再执行bin/up和拉取完整版的texlive，因为bin/up实际上执行了docker run，所以bin/up后再改配置文件就于事无补了。

\item bin/up实际上做了个docker run，bin/start是docker start。
修改variables.env或者overleaf.rc后，要重新使用bin/up生成新容器才能生效，同时也意味着你对旧容器做的更改会丢失（包括完整安装的textlive等），所以最好一开始就把所有配置项都调整好，不要后期重新bin/up。project应该是存在mongodb，不会丢失，但是完整的texlive会丢。

\item Overleaf 映像下载完成后会自动创建并启动容器，我们先按下 Ctrl + C 停止容器，或直接叉掉当前窗口，后面会实现后台运行。
\end{itemize}



4、接下来让 Overleaf 容器在后台运行：\\
\verb|bin/start|

5、Overleaf 容器启动之后，可以打开 http://localhost/launchpad 注册管理员帐户。之后我们就可以用这个帐户登录 Overleaf 平台。



\subsubsection{docker-compose.base.yml}
% Define colors for syntax highlighting
\definecolor{codegreen}{rgb}{0,0.6,0}
\definecolor{codegray}{rgb}{0.5,0.5,0.5}
\definecolor{codepurple}{rgb}{0.58,0,0.82}
\definecolor{backcolour}{rgb}{0.95,0.95,0.92}

% Configure listings package
\lstdefinestyle{mystyle}{
    backgroundcolor=\color{backcolour},   
    commentstyle=\color{codegreen},
    keywordstyle=\color{magenta},
    numberstyle=\tiny\color{codegray},
    stringstyle=\color{codepurple},
    basicstyle=\ttfamily\footnotesize,
    breakatwhitespace=false,         
    breaklines=true,                 
    captionpos=b,                    
    keepspaces=true,                 
    numbers=left,                    
    numbersep=5pt,                  
    showspaces=false,                
    showstringspaces=false,
    showtabs=false,                  
    tabsize=2
}

\lstset{style=mystyle}

% Define YAML language for listings
\lstdefinelanguage{YAML}{
  basicstyle=\ttfamily\small,
  keywordstyle=\bfseries\ttfamily\color{blue!70!black},
  commentstyle=\itshape\ttfamily\color{red!60!black},
  identifierstyle=\ttfamily\color{black},
  stringstyle=\ttfamily\color{orange!50!black},
  morecomment=[l]{\#},
  morekeywords={true,false,null,yaml,on,off},
  literate =     {---}{{\textcolor{blue}{---}}}3 % literal dash-dash-dash as blue
             {-}{{-}}1           % replace - with itself but allow it to be colored
             {...}{{\ldots}}1
}



The following is the content of the \texttt{docker-compose.base.yml} file along with comments explaining each line:

\begin{lstlisting}[language=YAML,caption=Docker Compose Configuration with Comments]
services: # 定义了所有要启动的服务列表。

    sharelatex: # 这是一个名为 `sharelatex` 的服务。

        restart: always # 指定容器总是自动重启，即使它因为某些原因退出或崩溃。
        
        image: "${IMAGE}" # 指定了用于创建容器的 Docker 镜像名称。这里的 `${IMAGE}` 是一个环境变量，在运行 `docker-compose` 命令时会被替换为实际的镜像名称。
        
        container_name: sharelatex # 设置了容器的名称为 `sharelatex`，这样可以通过这个名字来管理容器。
        
        volumes:
            - "${OVERLEAF_DATA_PATH}:${OVERLEAF_IN_CONTAINER_DATA_PATH}" # 将主机上的目录挂载到容器内的指定路径。`${OVERLEAF_DATA_PATH}` 和 `${OVERLEAF_IN_CONTAINER_DATA_PATH}` 是环境变量，分别表示主机和容器内数据存储的路径。
            
        ports:
            - "${OVERLEAF_LISTEN_IP:-127.0.0.1}:${OVERLEAF_PORT:-80}:80" # 映射主机端口到容器端口。`${OVERLEAF_LISTEN_IP}` 和 `${OVERLEAF_PORT}` 是环境变量，默认值分别为 `127.0.0.1` 和 `80`。这意味着默认情况下，ShareLaTeX 会在本地回环地址的 80 端口上监听，并将请求转发到容器内的 80 端口。
            
        environment:
          GIT_BRIDGE_ENABLED: "${GIT_BRIDGE_ENABLED}" # 设置了一系列环境变量，这些变量在容器内部可以被应用程序读取。其中 `${GIT_BRIDGE_ENABLED}`, `${REDIS_HOST}`, 和 `${REDIS_PORT}` 是从外部传入的环境变量；
          GIT_BRIDGE_HOST: "git-bridge" # 而 `GIT_BRIDGE_HOST`, `GIT_BRIDGE_PORT`, 和 `V1_HISTORY_URL` 则是固定的值。
          GIT_BRIDGE_PORT: "8000"
          REDIS_HOST: "${REDIS_HOST}"
          REDIS_PORT: "${REDIS_PORT}"
          V1_HISTORY_URL: "http://sharelatex:3100/api"
          
        env_file:
            - ../config/variables.env # 指定了额外的环境变量文件，位于项目根目录下的 `../config/variables.env` 文件中的键值对也会被加载到容器环境中。
            
        stop_grace_period: 60s # 容器停止前等待的时间长度。如果在此期间容器还没有停止，则强制终止它。这里设置的是 60 秒。
\end{lstlisting}







\subsection{安装宏包}
Overleaf 映像自带的宏包很少，为了编译我们的 LaTeX 文档，往往需要安装其他宏包。

1、在容器运行的状态下进入容器 Shell：\\
\verb|bin/shell|

2、更新并安装全部宏包：
\begin{verbatim}
# 使用中科大镜像源
tlmgr repository set https://mirrors.ustc.edu.cn/CTAN/systems/texlive/tlnet
# 更新已安装的全部宏包，包括 tlmgr 自身
tlmgr update --self --all
# 安装 scheme-full 套件（CTAN 上的全部宏包）
tlmgr install scheme-full
\end{verbatim}
注意: 安装全部宏包的时间大约要 40 分钟左右。

3、安装依赖包：
{\tiny
\begin{verbatim}
apt update
apt install ttf-mscorefonts-installer  # 安装 Windows 字体
apt install inkscape python3-pygments  # 安装 TikZ & Syntax highlighter 依赖

# 导入 TeX Live 字体
ln -s /usr/local/texlive/2024/texmf-var/fonts/conf/texlive-fontconfig.conf /etc/fonts/conf.d/09-texlive.conf
fc-cache -fsv  # 刷新字体缓存
\end{verbatim}
}


\subsection{配置Latex Windows或自定义字体}
Refs:\url{https://blog.xial.moe/index.php/2023/09/14/bushuzijideoverleaffuwu}

1、将Windows字体库（即目录 \verb|C:\windows\fonts|）上传到 host 机

2、在 host 机下把 fonts 目录打包并传到 sharelatex 容器中
\begin{verbatim}
# 进入fonts目录
$ cd fonts/
# 删除其中的.fon字体文件（该种格式文件在后面建立字体目录时会报错），
只保留TrueType和OpenType字体，即.ttf和.otf
# 一般地，如果只需要其中特定的中文字体，只需要上传需要的字体即可
$ rm -r *.fon
# 返回上层目录并打包
$ cd ..
$ tar -zcvf winfonts.tar.gz fonts/
# 把压缩文件传到sharelatex容器的root目录下
$ docker cp winfonts.tar.gz sharelatex:/root
\end{verbatim}

3、在容器中安装Windows字体
\begin{verbatim}
# 进入容器的命令行界面
$ docker exec -it sharelatex bash
# 通过安装wqy字体同时安装xfont工具
$ apt-get install xfonts-wqy
# 进入root目录，解压winfonts.tar.gz，并剪切到系统字体目录下
$ cd ~
$ tar -zxvf winfonts.tar.gz
$ mv winfonts /usr/share/fonts/
# 进入字体目录安装字体
$ cd /usr/share/fonts/winfonts
$ mkfontscale
$ mkfontdir
$ fc-cache -fv
# 检查确认中文字体安装成功
$ fc-list :lang=zh-cn
\end{verbatim}

4、回到 ShareLaTeX 网站，创建一个新项目，使用 CTEX 宏集和 XeLaTeX 编译器，即可生成中文 pdf。（详见 CTEX 宏集手册）

5、因为学校论文的需要，有时需要其他中文字体，例如Adobe宋体，其安装流程与上述基本一致。（PS：在GitHub上可以找到公开的Adobe Song Std字体，涉及版权这里不放链接）


PS：你可能还要安装一些中文字体，可以通过下面的命令将你的字体目录 fonts 拷贝到 Overleaf 容器的字体目录：\\
\verb|docker cp fonts sharelatex:/usr/local/share|



\subsection{备份数据或迁移}
参见官方文档。
\url{https://github.com/overleaf/overleaf/wiki/Data-and-Backups#Folders-in-detail}



\section{Troubleshooting}
\subsection{无法安装宏包}
在安装宏包时你可能遇到如下错误：
\begin{verbatim}
tlmgr: Local TeX Live (2023) is older than remote repository (2024).
Cross release updates are only supported with
  update-tlmgr-latest(.sh/.exe) --update
See https://tug.org/texlive/upgrade.html for details.
\end{verbatim}
这是映像使用的 TeX Live 版本过低导致的。如果你当前使用的不是最新的 ShareLaTeX 映像，可以使用下面的命令更新映像：\\
\verb|bin/upgrade|\\
如果你已经在使用最新的 ShareLaTeX 映像，那说明官方还没有给最新的 TeX Live 打映像。你可以设置 TeX Live 使用旧版镜像源，或者更新当前容器中的 TeX Live 版本，也可以自己打一个包含新版 TeX Live 的映像。

设置旧版镜像源：
{\tiny
\begin{verbatim}
tlmgr option repository https://mirrors.tuna.tsinghua.edu.cn/tex-historic-archive/systems/texlive/2021/tlnet-final/
\end{verbatim}
}



\subsection{打开 Web 界面时遇到 502 错误}
如果遇到 502 错误，可能是端口冲突导致的。进入 config/overleaf.rc 修改一下端口号：\\
\verb|OVERLEAF_PORT=8080|



\subsection{部署到服务器上之后无法访问}
Overleaf 的默认监听端口是 127.0.0.1:80，也就是只监听来自本地的 HTTP 访问请求。如果要开放所有 IP 的访问请求，需要修改 config/overleaf.rc 设置：
\begin{verbatim}
OVERLEAF_LISTEN_IP=0.0.0.0
OVERLEAF_PORT=80
\end{verbatim}
这样所有地址都可以访问服务器的 80 端口来获取 Overleaf 服务。



\subsection{无法启动 ShareLaTeX 容器}
如果你在运行 bin/up 的过程中遇到这样的错误提示：

Initiating Mongo replica set...
Rebranding from ShareLaTeX to Overleaf
  Starting with Overleaf CE and Server Pro version 5.0.0 the environment variables will use the Overleaf brand.
  Previous versions used the ShareLaTeX brand for environment variables.
  Your config/variables.env has 60 entries matching 'OVERLEAF\_', expected prefix 'SHARELATEX\_'.
  Please align your config/version with the naming scheme of variables in config/variables.env.
说明当前配置设置的 Overleaf 版本过低。我们需要使用大版本号为 5 以上的版本。

编辑 config/version，使用新的版本号替代旧的版本号。你可以在 Docker Hub 上查看最新的版本号：\\
\url{https://hub.docker.com/r/sharelatex/sharelatex/tags}



\subsection{无法启动 Mongo 容器}
运行 bin/up 命令时遇到如下错误：\\
MongoNetworkError: connect ECONNREFUSED 127.0.0.1:27017\\
Waiting for Mongo...\\
Please configure ALL\_TEX\_LIVE\_DOCKER\_IMAGES in /home/xiao/toolkit/config/variables.env\\
You can read more about configuring Sandboxed Compiles in our documentation:\\
  https://github.com/overleaf/overleaf/wiki/Server-Pro:-sandboxed-compiles\#changing-the-texlive-image

这可能是由于 Mongo 访问绑定挂载目录的 bug 导致的。要解决这个问题，我们需要将 Mongo 和 Redis 的数据存储位置由绑定挂载目录改为命名卷。方法为：创建 config/docker-compose.override.yml 文件，填入以下内容：
\begin{verbatim}
# 创建卷
volumes:
  mongo-data:
  redis-data:

# 挂载卷
services:
  mongo:
    volumes:
      - mongo-data:/data/db

  redis:
    volumes:
      - redis-data:/data
\end{verbatim}
然后重新运行 bin/up。




































